% Generated draft: requires mathematical and source review.
\documentclass[11pt]{article}
\usepackage[margin=25mm]{geometry}
\usepackage{fontspec}
\setmainfont{Latin Modern Roman}
\usepackage{amsmath,amssymb,mathtools,unicode-math}
\setmathfont{Latin Modern Math}
\usepackage{enumitem,needspace,xurl,hyperref}
\setlength{\parindent}{0pt}
\setlength{\parskip}{5pt}
\setlength{\emergencystretch}{3em}
\newcommand{\N}{\mathbb{N}}
\newcommand{\Z}{\mathbb{Z}}
\newcommand{\Q}{\mathbb{Q}}
\newcommand{\R}{\mathbb{R}}
\newcommand{\C}{\mathbb{C}}
\newcommand{\F}{\mathbb{F}}
\newcommand{\A}{\mathbb{A}}
\newcommand{\PP}{\mathbb{P}}
\newcommand{\E}{\mathbb{E}}
\newcommand{\eps}{\varepsilon}
\newcommand{\dd}{\,\mathrm{d}}
\begin{document}
\section*{UnsolvedMath ID: 30003080}
\textbf{Source problem code:} catalogue entry
\subsection{Definitions and mathematical statement}
The supplied record asks whether a combinatorial invariance property of the multiplicity index implies Terao's conjecture for line arrangements. It does not define the multiplicity index, specify the ground field or ambient space of the point sets, or state the range of the integer parameter \(k\). Consequently, the following is a conditional formulation with explicitly unresolved parameters, rather than a fully specified conjecture.

Let \(I\) denote the intended range of integers \(k\), let \(\mathcal P_k\) denote the intended class of sets of \(2k+1\) points, and let \(\sim_{\mathrm{comb}}\) denote the intended equivalence relation expressing equality of combinatorial data. Write \(m_Z\) for the source's multiplicity index of \(Z\); its definition must be recovered from the cited source before this formulation can be made complete. The proposed combinatorial property is
\[
\mathsf C:\qquad
\forall k\in I\;\forall Z,Z'\in\mathcal P_k,\qquad
Z\sim_{\mathrm{comb}}Z'
\ \Longrightarrow\
\bigl(m_Z=k\ \Longleftrightarrow\ m_{Z'}=k\bigr).
\]
This asserts invariance of the predicate \(m_Z=k\), not necessarily invariance of the entire numerical invariant \(m_Z\).

For clarity, the usual projective-plane formulation of the proposed conclusion can be described with an unspecified ground field \(K\). For a finite arrangement \(\mathcal A\) of distinct lines in \(\mathbb P^2_K\), choose homogeneous linear forms \(\alpha_H\in S=K[x_0,x_1,x_2]\) defining its lines. Its central cone is the arrangement of planes \(\ker(\alpha_H)\subseteq K^3\). Its intersection lattice \(L(\mathcal A)\) consists of all intersections of these planes, including the empty intersection \(K^3\), ordered by reverse inclusion. Define the module of logarithmic derivations by
\[
D(\mathcal A)=
\left\{\theta\in\operatorname{Der}_K(S):
\theta(\alpha_H)\in\alpha_H S\text{ for every }H\in\mathcal A\right\}.
\]
Call \(\mathcal A\) free if \(D(\mathcal A)\) is a free \(S\)-module. In this interpretation, Terao's conjecture is the assertion
\[
\mathsf T_K:\qquad
\forall\mathcal A,\mathcal B,\qquad
L(\mathcal A)\cong L(\mathcal B)
\ \Longrightarrow\
\bigl(\mathcal A\text{ is free}\ \Longleftrightarrow\
\mathcal B\text{ is free}\bigr),
\]
where both arrangements are over the same field \(K\).

The source question is whether
\[
\mathsf C\ \Longrightarrow\ \mathsf T_K.
\]
The projective-plane interpretation above explains a standard meaning of the conclusion; the record does not establish that this is its exact intended field and domain. It also does not specify whether the combinatorial equivalence of point sets is defined through the intersection lattices of their dual line arrangements. These missing choices, and the definition of \(m_Z\), remain unresolved.

The supplied dated curation describes the implication as established while describing the required combinatorial property as unknown in general. This is a supplied assessment, not a mathematical verification made in this draft.
\subsection{Short English statement}
Does invariance of the condition \(m_Z=k\) under the intended combinatorial equivalence, for every intended set of \(2k+1\) points, imply that freeness of line arrangements is determined by their intersection lattices? The definition of \(m_Z\), the ground field, the point-set domain, and the range of \(k\) are missing. The supplied curation labels the problem open but says the implication is established and its premise remains unknown in general; that assessment requires independent review.
\subsection{Sources}
\begin{itemize}
\item Supplied local UnsolvedMath registry record, permanent ID \(30003080\), title ``Combinatorial Multiplicity Indices and Terao Conjecture,'' source problem code \texttt{catalogue entry}, in the Oberwolfach Reports Open Problems collection. The project identifies \texttt{data/unsolvedmath/problems.json} and the supplementary \texttt{catalog/problems.json} as local registry resources. No individual proposer is supplied; the proposed year is 2016.
\item Supplied report citation: \emph{Mini-Workshop: Arrangements of Subvarieties, and their Applications in Algebraic Geometry} (2016), Oberwolfach Reports, pp.~653--702, DOI \texttt{10.4171/owr/2016/14}. \url{https://doi.org/10.4171/owr/2016/14}. The background describes Oberwolfach Reports as documenting mathematical workshops and their open-problem discussions.
\item Supplied literature source: Cook--Harbourne--Migliore--Nagel, \url{https://arxiv.org/abs/1602.02300}. The record supplies no title or pinpoint reference for the relevant claim.
\item Supplied literature source: Pagaria, \url{https://doi.org/10.1112/jlms.12829}. The record supplies no title or pinpoint reference for the relevant claim.
\item Literature assessment dated exactly \texttt{2026-08-21}: the supplied curation labels the status ``open'' and states that the implication to Terao's conjecture is established, while the required combinatorial multiplicity-index property is not known in general. The background calls this dated literature triage derived from the supplied OWR curation and requests independent verification of time-sensitive status and cited claims. This is supplied dated curation, not independently verified here.
\item The registry labels the statement \texttt{corrected\_verified}, with the explanation that notation, typography, and sentence structure were repaired without changing mathematical content. This supplied editorial label does not establish mathematical correctness.
\end{itemize}
\end{document}
